\documentclass[10pt,a4paper]{amsart}
\usepackage[margin=19mm]{geometry}
\usepackage{amsmath,amssymb,mathtools}
\usepackage[scr=rsfs,cal=euler]{mathalfa}
\usepackage{xfrac}
\usepackage[unicode,hidelinks,bookmarks=false]{hyperref}
\newcommand{\ie}{i.e.\ }
\newcommand{\cf}{cf.\ }
\newcommand{\resp}{respectively\ }
\newcommand{\Subsection}{\subsection}
\providecommand{\nobreakdash}{\nobreak\hbox{-}}
\setlength{\emergencystretch}{2em}
\multlinegap0pt
\usepackage{bm}
\usepackage{bbm}
\usepackage{dsfont}
\usepackage{xspace}
\usepackage{cancel}
\usepackage{mathtools}
\usepackage{amsthm}
\makeatletter
\@ifundefined{theorem}{\newtheorem{theorem}{\bf Theorem}}{}
\@ifundefined{corollary}{\newtheorem{corollary}[theorem]{\bf Corollary}}{}
\@ifundefined{proposition}{\newtheorem{proposition}[theorem]{\bf Proposition}}{}
\makeatother
\title[Positive Solutions for Sublinear Equations with Compact Posi]{Positive Solutions for Sublinear Equations with Compact Positivity-Improving Resolvent}
\author{Tomasz Klimsiak}
\date{Source: arXiv:2605.10570v2 (CC BY 4.0)}
\begin{document}
\maketitle

\noindent\emph{Source and licence.} Positive Solutions for Sublinear Equations with Compact Positivity-Improving Resolvent. Version arXiv:2605.10570v2 (CC BY 4.0), distributed under the Creative Commons Attribution 4.0 International licence, as recorded in the arXiv licence metadata for this exact version and shown on the abstract page https://arxiv.org/abs/2605.10570v2. Full rights notice: licenses/arxiv-2605.10570-CC-BY-4.0.txt.\par\smallskip

\noindent\textit{Editorial scope.} Counted unit: the Corollary whose source label is \texttt{cor.weakerF3} in the article (source0.tex lines 2420--2444, source label \texttt{cor.weakerF3}), delivered together with the article's own complete original proof of it (source0.tex lines 2446--2813, one \texttt{proof} environment of 368 lines). No mathematics is written, repaired or re-derived: statement and proof are the article's own text line for line. Required context retained uncounted: eq1.1 (lines 198--200); lm.upseq (lines 1013--1018); lm2 (lines 1174--1182); prop.shifted-equation (lines 1755--1774); th.mainm (lines 1809--1816); th.main (lines 2257--2265). Every definition, display and label of those lines is retained unchanged. Omitted from this package and not redistributed: 1--197, 201--1012, 1019--1173, 1183--1754, 1775--1808, 1817--2256, 2266--2419, 2445--2445, 2814--3476. Nothing inside a retained excerpt is omitted or altered: every retained body is delivered exactly as the article's lines are written, so no omission receipt of any kind accompanies this package. Citations: the retained excerpts carry no citation command at all, so no bibliography entry of the article is transcribed and no reference is redistributed. Reference adaptation: none was needed and none was made; every reference inside the delivered unit resolves inside it, so no unresolved reference or citation is produced. Printed slips kept literal: no printed word, symbol, hypothesis or number of the counted statement or of its proof was altered. Layout and class: the article's own class is \texttt{amsart} with packages including fontenc, inputenc, lmodern, microtype, geometry, mathrsfs, amsmath, amssymb, amsthm, enumerate, hyperref; the wrapper is the lane's pinned \texttt{amsart} editorial wrapper at 10pt on A4, which restates the theorem-like environments the retained text needs and adds the article's own macro definitions that the retained text needs (none), with extra wrapper packages (bm, bbm, dsfont, xspace, cancel, mathtools). The delivered numbering is the unit's own. Assets: no retained span contains a figure environment, a graphics-inclusion command, a PDF-page inclusion, a table or a TikZ picture; no image file is copied into this package, the package declares no asset dependency and no image file is redistributed with it. Licence: version arXiv:2605.10570v2 of this article is distributed under the Creative Commons Attribution 4.0 International licence, as recorded in the arXiv licence metadata for this exact version and shown on the abstract page https://arxiv.org/abs/2605.10570v2. A full rights notice is delivered at licenses/arxiv-2605.10570-CC-BY-4.0.txt. Characters: every retained excerpt is pure ASCII, so the delivered engine, XeLaTeX with its default Latin Modern text font, reports no missing character.
\begin{equation}\label{eq1.1}
    -Lu=f(x,u)\quad\text{in }E,
\end{equation}

\begin{theorem}
\label{lm.upseq}
Assume that $(T_t)$ is sub-Markovian. For any sequence $(v_n)$ of
supermedian functions there exists a subsequence $(v_{n_k})$ which
converges $m$-a.e. in $E$, possibly to $+\infty$.
\end{theorem}

\begin{proposition}
\label{lm2}
Assume that $(T_t)$ is sub-Markovian.   Suppose that  $v,w$ are  supermedian functions.
Set $u:=Rf-w$  for some $f\in L^p(E;m)$.
Then
\[
(u-v)^+\le R(\mathbf1_{\{u>v\}}f).
\]
\end{proposition}

\begin{proposition}
\label{prop.shifted-equation}
Let $\gamma$ be a bounded supermedian function,
$\kappa\geq 0$, 
\[
L_\kappa:=L-\kappa I,
\qquad
f^\kappa(x,y):=f(x,y)+\kappa y,
\]
and    $u\in L^p(E;m)$ be a positive function. Then  $u$ is a solution to
\begin{equation}
\label{eq.res123}
u=\gamma+Rf(\cdot,u)
\end{equation}
if and only if it is a solution to
\begin{equation}
\label{eq.res1234}
u=\gamma-\kappa R_\kappa\gamma+R_\kappa f^\kappa(\cdot,u).
\end{equation}
\end{proposition}

\begin{theorem}
\label{th.mainm}
Assume that $(T_t)$ is sub-Markovian and (F1)-(F2) and (F3)$_1$. If 
\[
\lambda_1(a_0)<0< \lambda_1(a_\infty),
\]
then there exists a strictly positive   solution to \eqref{eq1.1}.
\end{theorem}

\begin{theorem}
\label{th.main}
Assume that (F1)-(F3) hold. If 
\begin{equation}
\label{eq.nine}
\lambda_1(a_0)<0< \lambda_1(a_\infty),
\end{equation}
then there exists a  strictly positive solution to \eqref{eq1.1}.
\end{theorem}

\begin{corollary}
\label{cor.weakerF3}
Assume that \textup{(F1)--(F2)} hold and suppose, in addition, that the following  condition \textup{(F3')} holds:
\begin{enumerate}
\item[{\rm (F3')}] there exists a strictly positive supermedian function $\psi$ such that:
\begin{enumerate}
\item[(i)] there exist $\delta>0$ and $C_\delta\ge0$ such that
\[
f^-(x,y)\le C_\delta y,
\qquad x\in E,\quad 0\le y\le \psi(x)\delta;
\]
\item[(ii)] for every $M>0$, the function 
\[
F_M(x):=\sup_{0\le y\le \psi(x) M}f^-(x,y)
\]
belongs to $L^p(E;m)$.
\end{enumerate} 
\end{enumerate}
If
\begin{equation}
\label{eq.weakerF3.spectral}
\lambda_1(a_0)<0<\lambda_1(a_\infty),
\end{equation}
then \eqref{eq1.1} admits a strictly positive solution.
\end{corollary}

\begin{proof}
As in the proof of Theorem \ref{th.main} we may reduce to the case where $(T_t)$ is sub-Markovian and
$\psi\equiv 1$ in (F3'). 
By Proposition \ref{prop.shifted-equation}, we may assume without
loss of generality that
\begin{equation}
\label{eq.localpos}
R1\le 1,\quad \text{and}\quad f(x,y)\ge0,
\qquad x\in E,\quad 0\le y\le\delta.
\end{equation}
Choose $l\ge1$ so large that
\begin{equation}
\label{eq.choose.l}
\lambda_1(a_0^l)<0<\lambda_1(a_\infty^l).
\end{equation}
For $k,n\ge l$, define, for $x\in E$, $y\ge 0$,
\begin{equation}
\label{eq.newfkn}
f_k(x,y)
:=
\min\{f^+(x,y),ky,k\}-f^-(x,y),\quad  f_{k,n}(x,y)
:=f^+_k(x,y)-\min\{f^-(x,y),ny\}.
\end{equation}
Then
\begin{equation}
\label{eq.estmt}
f_{k,n}^+(x,y)\le ky\wedge k,
\qquad
f_{k,n}^-(x,y)\le ny.
\end{equation}
In particular, $f_{k,n}$ satisfies \textup{(F1)--(F3)}.
Moreover, by \eqref{eq.localpos},
\[
a_0(f_{k,n})=a_0\wedge k=a_0^k,\quad a_\infty(f_{k,n})=\bigl(a_\infty\vee(-n)\bigr)\wedge0.
\]
Consequently, for $k,n\ge l$,
\[
\lambda_1(a_0(f_{k,n}))
=
\lambda_1(a_0^k)
\le \lambda_1(a_0^l)<0,
\]
while, by monotonicity of the generalized principal eigenvalue,
\[
\lambda_1(a_\infty(f_{k,n}))
\ge
\lambda_1(a_\infty^n)
\ge
\lambda_1(a_\infty^l)>0.
\]
Thus $f_{k,n}\in\mathscr F_1$.
Set
\[
u_{k,n}:=\mathcal S_1(f_{k,n}),
\qquad k,n\ge l.
\]
Since $f_{k,n}$ is decreasing in $n$ and increasing in $k$, the
monotonicity of $\mathcal S_1$ gives
\begin{equation}
\label{eq.newmon}
u_{k,n}\ge u_{k,n+1},
\qquad
u_{k,n}\le u_{k+1,n},
\qquad k,n\ge l.
\end{equation}

{\bf Step 1. Letting $n\to\infty$.}  For fixed $k$, define
\[
u_k:=\inf_{n\ge l}u_{k,n}.
\]
Since
\[
0\le u_{k,n}\le Rk\le k,
\]
we have
\begin{equation}
\label{eq.untoUK}
u_{k,n}\longrightarrow u_k
\quad\text{a.e. and in }L^p(E;m).
\end{equation}
We first show that $u_k$ is nontrivial. Suppose, to the contrary,
that $u_k=0$. Then
\[
c_n:=\|u_{k,n}\|\searrow0.
\]
Set $w_{k,n}:=u_{k,n}/c_n$. Then 
\[
w_{k,n}
=
R\left(f_{k,n}(\cdot,u_{k,n})/c_n
\right).
\]
Since $f_{k,n}^-(\cdot,u_{k,n})\le n u_{k,n}$ and
$f_{k,n}^+(\cdot,u_{k,n})\le k u_{k,n} $, both parts belong to
$L^p(E;m)$.
Put
\[
\eta_n
:=
R\left(
f_{k,n}^+(\cdot,u_{k,n})/c_n
\right),
\qquad
\zeta_n
:=
R\left(
f_{k,n}^-(\cdot,u_{k,n})/c_n
\right).
\]
Then
\[
0\le \eta_n\le kRw_{k,n},
\qquad
w_{k,n}=\eta_n-\zeta_n.
\]
Since $\|w_{k,n}\|=1$, the sequences $(\eta_n)$ and $(\zeta_n)$ are
bounded in $L^p(E;m)$. Both consist of supermedian functions.
Hence, by Theorem \ref{lm.upseq}, after passing to a subsequence,
$\eta_n$ and $\zeta_n$ converge a.e. Consequently,
\[
w_{k,n}\longrightarrow w_k
\qquad\text{a.e.}
\]
along the same subsequence.
Furthermore,
\[
0\le w_{k,n}\le \eta_n\le kRw_{k,n}.
\]
By compactness of $R$, after passing to a further subsequence,
$(Rw_{k,n})$ converges in $L^p(E;m)$. It follows that
$(w_{k,n}^p)$ is uniformly integrable. Therefore, by Vitali's
theorem,
\begin{equation}
\label{eq.wnstrong}
w_{k,n}\longrightarrow w_k
\quad\text{in }L^p(E;m).
\end{equation}
In particular, $\|w_k\|=1$.
Let $\varphi_l'$ be a dual principal eigenfunction associated with
$-L-a_0^l$, and choose $\theta>0$ so large that
\[
a_0^l+\lambda_1(a_0^l)+\theta\ge0
\qquad\text{on }E.
\]
By Proposition \ref{prop.shifted-equation},
\[
w_{k,n}
=
R_\theta \left [(f_{k,n}(\cdot,u_{k,n})+\theta u_{k,n})/c_n\right].
\]
Fix $r>0$. By Proposition \ref{lm2}
\[
(w_{k,n}-r)^+
\le
R_\theta
\left(
\mathbf1_{\{w_{k,n}>r\}}(f_{k,n}(\cdot,u_{k,n})+\theta u_{k,n})/c_n
\right).
\]
Consequently,
\begin{equation}
\label{eq.truncwn}
w_{k,n}\wedge r
\ge
R_\theta
\left(
\mathbf1_{\{w_{k,n}\le r\}}(f_{k,n}(\cdot,u_{k,n})+\theta u_{k,n})/c_n
\right).
\end{equation}
Pairing  \eqref{eq.truncwn} with the non-negative function $\bigl(a_0^l+\lambda_1(a_0^l)+\theta\bigr)\varphi_l'$ gives
\begin{equation}
\label{eq.truncpair}
\left\langle w_{k,n}\wedge r,\bigl(a_0^l+\lambda_1(a_0^l)+\theta\bigr)\varphi_l'\right\rangle
\ge\left\langle\mathbf1_{\{w_{k,n}\le r\}}w_{k,n}\left(f_{k,n}(\cdot,u_{k,n})/u_{k,n}+\theta\right),\varphi_l'\right\rangle.
\end{equation}
Since, under the contradiction assumption, $u_{k,n}\to0$ a.e. as $n\to \infty$, we have,  by the definition of $a_0$,
\[
\liminf_{n\to\infty}
\frac{f_{k,n}(x,u_{k,n}(x))}{u_{k,n}(x)}
\ge
a_0(x)\wedge k
\ge a_0^l(x).
\]
Observe that
\[
\{w_{k,n}\le r\}
=
\{u_{k,n}\le r c_n\}.
\]
Since $c_n\to0$, for every fixed $r$ we have
$r c_n\le\delta$ for all sufficiently large $n$. Hence, by
\eqref{eq.localpos}, the integrand on the right-hand side of
\eqref{eq.truncpair} is nonnegative for all sufficiently large $n$.
Choose $r>0$ such that
\[
m(\{w_k=r\})=0.
\]
Using \eqref{eq.wnstrong}, Fatou's lemma, and
\eqref{eq.truncpair}, we obtain
\begin{equation}
\label{eq.limit.r}
\left\langle
w_k\wedge r,
\bigl(a_0^l+\lambda_1(a_0^l)+\theta\bigr)\varphi_l'
\right\rangle\ge\left\langle\mathbf1_{\{w_k<r\}}w_k(a_0^l+\theta),\varphi_l'\right\rangle.
\end{equation}
Taking a sequence $r_j\uparrow\infty$ such that
$m(\{w_k=r_j\})=0$ and letting $j\to\infty$ in
\eqref{eq.limit.r}, we get
\[
\left\langle
w_k,
\bigl(a_0^l+\lambda_1(a_0^l)+\theta\bigr)\varphi_l'
\right\rangle
\ge
\left\langle
w_k(a_0^l+\theta),
\varphi_l'
\right\rangle .
\]
Thus
\[
\lambda_1(a_0^l)
\langle w_k,\varphi_l'\rangle
\ge0,
\]
a contradiction.
Hence $u_k$ is nontrivial.  
We next pass to the limit as $n\to\infty$ in the equation for
$u_{k,n}$. From \eqref{eq.estmt}--\eqref{eq.untoUK}
and the dominated convergence theorem we infer that
\begin{equation}
\label{eq.posnlim}
f_{k,n}^+(\cdot,u_{k,n})=f_k^+(\cdot,u_{k,n})
\longrightarrow
f_k^+(\cdot,u_k)
\quad\text{in }L^p(E;m)
\end{equation}
(see \eqref{eq.newfkn}). It remains to treat the negative part. Fix $M>0$ such that $m(\{u_k=M\})=0$. On
$\{u_{k,n}\le M\}$,
\[
f_{k,n}^-(\cdot,u_{k,n})
\le
F_M,
\]
and \textup{(F3')(ii)} together with dominated convergence yields
\begin{equation}
\label{eq.localq}
\mathbf1_{\{u_{k,n}\le M\}}f_{k,n}^-(\cdot,u_{k,n})
\longrightarrow
\mathbf1_{\{u_k\le M\}}f^-(\cdot,u_k)
\quad\text{in }L^p(E;m).
\end{equation}
On the other hand, Proposition \ref{lm2} gives
\[
(u_{k,n}-M)^+
+
R\left(
\mathbf1_{\{u_{k,n}>M\}}f_{k,n}^-(\cdot,u_{k,n})
\right)
\le
R\left(
\mathbf1_{\{u_{k,n}>M\}}
f_k^+(\cdot,u_{k,n})
\right).
\]
Since $f_k^+(\cdot,u_{k,n})\le k u_{k,n}$ and
$u_{k,n}\le u_{k,l}$,
\begin{equation}
\label{eq.tailn}
R\left(
\mathbf1_{\{u_{k,n}>M\}}f_{k,n}^-(\cdot,u_{k,n})
\right)
\le
kR\left(
\mathbf1_{\{u_{k,l}>M\}}u_{k,l}
\right).
\end{equation}
The right-hand side converges to zero in $L^p(E;m)$ as
$M\to\infty$, uniformly in $n$.
Using the kernel representation of $R$ and Fatou's lemma, we also
obtain
\[
R\left(
\mathbf1_{\{u_k>M\}}f^-(\cdot,u_k)
\right)
\le
kR\left(
\mathbf1_{\{u_{k,l}>M\}}u_{k,l}
\right).
\]
Combining this estimate with \eqref{eq.localq} and
\eqref{eq.tailn}, first letting $n\to\infty$ and then
$M\to\infty$, yields
\begin{equation}
\label{eq.negpotlim}
Rf_{k,n}^-(\cdot,u_{k,n})
\longrightarrow
Rf^-(\cdot,u_k)
\quad\text{in }L^p(E;m).
\end{equation}
Passing to the limit in
\[
u_{k,n}+Rf_{k,n}^-(\cdot,u_{k,n})
=
Rf_k^+(\cdot,u_{k,n})
\]
and using \eqref{eq.untoUK}, \eqref{eq.posnlim}, and
\eqref{eq.negpotlim}, we obtain
\begin{equation}
\label{eq.uk}
u_k+Rf^-(\cdot,u_k)=Rf_k^+(\cdot,u_k).
\end{equation}
We now show that $u_k$ is strictly positive. Since
$f_k(x,y)\ge0$ for $0\le y\le\delta$, the localized concave
truncation argument with the function $\rho_n$ used in Step~2 of the proof of
Theorem \ref{th.mainm} 
applies to \eqref{eq.uk} and shows that $u_k\gg0$.


{\bf Step 2. Letting $k\to\infty$.}  By \eqref{eq.newmon},
\[
u_k\le u_{k+1},
\qquad k\ge l.
\]
Suppose first that $\sup_{k\ge l}\|u_k\|<\infty$. Then
\[
u_k\uparrow u
\quad\text{a.e. and in }L^p(E;m)
\]
for some $u\in L^p(E;m)$. Since every $u_k$ is strictly positive,
$u\gg0$. By \textup{(F1)},
\[
f_k^+(\cdot,u_k)\le f^+(\cdot,u_k)\le h+\lambda u,
\]
and therefore
\[
f_k^+(\cdot,u_k)\longrightarrow f^+(\cdot,u)\quad\text{in }L^p(E;m).
\]
For the negative part we argue as above. For every $M>0$ such that
$m(\{u=M\})=0$, \textup{(F3')(ii)} gives
\[
\mathbf1_{\{u_k\le M\}}f^-(\cdot,u_k)\longrightarrow\mathbf1_{\{u\le M\}}f^-(\cdot,u)\quad\text{in }L^p(E;m).
\]
Moreover, Proposition \ref{lm2} and \textup{(F1)} imply
\[
R\left(\mathbf1_{\{u_k>M\}}f^-(\cdot,u_k)\right)\le 
R\left(\mathbf1_{\{u_k>M\}}f_k^+(\cdot,u_k)\right)\le R\left(\mathbf1_{\{u>M\}}(h+\lambda u)\right).
\]
The last term converges to zero in $L^p(E;m)$ as $M\to\infty$.
Using again the kernel representation and Fatou's lemma for the
limit function, we conclude that
\[
Rf^-(\cdot,u_k)\longrightarrow Rf^-(\cdot,u) \quad\text{in }L^p(E;m).
\]
Passing to the limit in \eqref{eq.uk} gives
\[
u+Rf^-(\cdot,u)=Rf^+(\cdot,u),
\]
so $u$ is a strictly positive solution of \eqref{eq1.1}.
It remains to exclude the possibility that
\[
\sup_{k\ge l}\|u_k\|=\infty.
\]
The argument is identical to that used in Step~3 of the proof of
Theorem~\ref{th.mainm}. Therefore the unbounded alternative cannot occur,
and the proof is complete.
\end{proof}

\end{document}
